\worksheettitle
  {Точки через равные промежутки}
  {\modulemark{M07} Версия учителя}

\begin{teacherbox}[Паспорт модуля]
  \textbf{Результат:} ученик различает позиции и промежутки, решает прямые и
  обратные задачи на равный шаг, находит конечный номер и переносит модель
  на ряд объектов с включёнными концами.

  \textbf{Предварительно:} измеряет длину и выполняет действия с натуральными
  числами.
\end{teacherbox}

\begin{teacherbox}[Методический акцент]
  Не начинайте с формулы. На коротком ряду ученик показывает каждый переход,
  после чего разность номеров получает смысл числа переходов.
\end{teacherbox}

\section{Вход и основная модель}

Во входе: $5$ точек и $4$ промежутка; промежутки находятся между соседними
точками.

\pointsvsintervalsfigure

\begin{practicebox}[1. Число промежутков]
  a) $6$; б) $19$; в) $1$; г) $15$.
\end{practicebox}

\studentcountstrip

\begin{practicebox}[2. По рисунку]
  $7-2=5$ промежутков; расстояние $5\cdot3=15$ см.
\end{practicebox}

\section{Прямая и обратная задачи}

\distanceexamplefigure

\begin{practicebox}[3. Расстояния]
  a) $(12-5)\cdot7=49$ мм;
  б) $(27-8)\cdot7=133$ мм;
  в) $(44-31)\cdot7=91$ мм.
\end{practicebox}

\reverseexamplefigure

В примере один промежуток равен $24:4=6$ см.

\begin{practicebox}[4. Конечный номер]
  a) $28:4=7$, $12+7=19$;
  б) $36:4=9$, $30-9=21$.
\end{practicebox}

\section{Ряды объектов}

\postrowfigure

\begin{practicebox}[5. Столбы]
  Семь столбов образуют $7-1=6$ промежутков; расстояние
  $6\cdot4=24$ м.
\end{practicebox}

\begin{practicebox}[6. Фонари]
  $54:6=9$ промежутков. При включённых крайних фонарях объектов на один
  больше: $9+1=10$ фонарей.
\end{practicebox}

\begin{teacherbox}[Граница правила]
  Прибавлять единицу можно только после проверки, что объект есть в обоих
  концах ряда. Если условие этого не сообщает, расположение концов нужно
  уточнить или рассмотреть случаи.
\end{teacherbox}

\section{Ошибка и контроль}

\begin{warningbox}[Ожидаемое исправление]
  Причина: точки приняты за промежутки. Между 1-й и 10-й точками
  \[
    10-1=9
  \]
  промежутков.
\end{warningbox}

\begin{checkpointbox}[Ответы]
  \begin{enumerate}[label=\textbf{\arabic*.}]
    \item $(19-7)\cdot5=60$ см;
    \item $10-4=6$ промежутков, $42:6=7$ мм;
    \item $32:4=8$ промежутков, $8+1=9$ столбов;
    \item каждый промежуток расположен между двумя соседними объектами;
      первый объект начинает ряд, а каждый следующий завершает новый
      промежуток.
  \end{enumerate}
\end{checkpointbox}

\begin{teacherbox}[Решение о маршруте]
  M07-R назначается при системном смешении точек и промежутков. M07-H
  достаточно при верной модели и единичной арифметической ошибке.
  M07\(\star\) предлагается после устойчивого базового контроля.
\end{teacherbox}
